\begin{apartats}

\apartat[1,25]{0,75}
Escriu les equacions paramètriques i contínua de la recta que passa per $P(1,-2,0)$ i és paral·lela a
la recta $r\colon\dfrac{x-3}{2}=y+1=\dfrac{z-2}{-3}$.

\begin{solucio}
Una recta paral·lela a $r$ té el mateix vector director, $\vec u=(2,1,-3)$.\\
Paramètriques: $\left\{\begin{aligned}x&=1+2t\\y&=-2+t\\z&=-3t\end{aligned}\right.$\qquad
Contínua: $\dfrac{x-1}{2}=\dfrac{y+2}{1}=\dfrac{z}{-3}$.
\end{solucio}

\begin{tria}{altres-formes}
\itemtria{implicita-parametrica}{1}{1,25}
Escriu en forma paramètrica la recta
$s\colon\left\{\begin{aligned}x+y-z&=1\\2x-y+z&=5\end{aligned}\right.$, i dona'n un punt i un vector
director.

\begin{solucio}
És un sistema compatible indeterminat. Amb $z=\lambda$: sumant les equacions, $3x=6$, $x=2$; i de la
primera, $y=1-x+z=\lambda-1$.
$s\colon\left\{\begin{aligned}x&=2\\y&=-1+\lambda\\z&=\lambda\end{aligned}\right.$ Un punt és
$(2,-1,0)$, i un vector director, $(0,1,1)$.
\end{solucio}

\itemtria{alineats-dos-parametres}{1}{1,25}
Troba els valors de $a$ i $b$ perquè els punts $A(1,0,2)$, $B(2,1,a)$ i $C(4,b,8)$ estiguin alineats.

\begin{solucio}
$\overrightarrow{AB}=(1,1,a-2)$ i $\overrightarrow{AC}=(3,b,6)$ han de ser proporcionals. La primera
component dona la raó 3: $b=3\cdot1=3$ i $6=3(a-2)$, és a dir, $a=4$.
\end{solucio}
\end{tria}

\begin{nomesllarg}
\apartat{0,75}
Escriu en forma implícita la recta $\dfrac{x-1}{2}=\dfrac{y+3}{-1}=\dfrac{z}{4}$.

\begin{solucio}
Igualant la primera fracció amb cadascuna de les altres: $-(x-1)=2(y+3)$ i $4(x-1)=2z$.
$\left\{\begin{aligned}x+2y+5&=0\\2x-z-2&=0\end{aligned}\right.$
\end{solucio}
\end{nomesllarg}

\end{apartats}
